\documentclass[11pt,a4paper]{article}
\usepackage[T2A]{fontenc}
\usepackage[utf8]{inputenc}
\usepackage[russian]{babel}
\usepackage{amsmath,amssymb,amsthm}
\usepackage[left=2cm,right=2cm,top=2cm,bottom=2cm]{geometry}
\usepackage{fancyvrb}
\usepackage[most]{tcolorbox}
\usepackage[unicode,hidelinks]{hyperref}
\setlength{\parindent}{0pt}\setlength{\parskip}{4pt}
\newtheorem{theorem}{Теорема}
\newtheorem{lemma}{Лемма}
\newtcolorbox{key}[1]{colback=blue!3,colframe=blue!45!black,boxrule=0.5pt,arc=1mm,title={#1},fonttitle=\bfseries\small}
\newtcolorbox{qa}{colback=gray!4,colframe=gray!55!black,boxrule=0.4pt,arc=1mm,breakable}

\title{\textbf{Задача 1.17: решение уравнения $1+e^{-2x^2}=\alpha^2e^{-2x}$\\ методом простой итерации}\\[4pt]\large Подробный разбор: математика и код}
\date{}
\begin{document}
\maketitle
\tableofcontents

\section{Что сделано и как запустить}
Файлы в папке \texttt{task117\_yarik/}:
\begin{itemize}
\item \texttt{iterative\_solver.cpp} — программа (устроена так же, как образец \texttt{iterative\_solver.cpp} из корня репозитория);
\item \texttt{report.pdf} (\texttt{report.tex}) — отчёт в формате образца \texttt{report.pdf};
\item \texttt{razbor.pdf} (\texttt{razbor.tex}) — этот разбор.
\end{itemize}
Сборка и запуск (нужен только компилятор C++, сторонние библиотеки не нужны):
\begin{Verbatim}[frame=single,fontsize=\small]
g++ -O2 -Wall -std=c++17 -o iterative_solver iterative_solver.cpp
./iterative_solver 2.5        # одно значение alpha (любое вещественное)
./iterative_solver -10 10     # все целые alpha от -10 до 10
\end{Verbatim}
Вывод — таблица из столбцов $\alpha$, $x$ (найденное решение), $n$ (номер приближения, удовлетворяющего точности $10^{-12}$).

\begin{key}{Суть решения в пяти строках}
\begin{enumerate}
\item При $\alpha=0$ решений нет (левая часть $\ge1$, правая $=0$).
\item При $\alpha\ne0$ логарифмируем: $2x+\ln(1+e^{-2x^2})=2\ln|\alpha|$.
\item Выражаем $x$: $x=G(x)=\ln|\alpha|-\tfrac12\ln(1+e^{-2x^2})$.
\item $|G'(x)|=\dfrac{2|x|}{e^{2x^2}+1}\le\dfrac12$ на всей прямой $\Rightarrow$ $G$ — сжатие с $k=\tfrac12$ $\Rightarrow$ (Банах) решение существует, единственно, итерации сходятся.
\item Останавливаемся, когда $\dfrac{|x_{n+1}-x_n|}{1-k}\le\varepsilon$, тогда $|x_n-x^*|\le\varepsilon$.
\end{enumerate}
\end{key}

\section{Анализ уравнения}
\subsection{Случай $\alpha=0$}
Левая часть $1+e^{-2x^2}>1$, а правая равна $0$. Решений нет. Программа в этом случае печатает \texttt{no solution}.

\subsection{Случай $\alpha\ne0$: переход к логарифмам}
Обе части уравнения строго положительны, поэтому логарифмирование — равносильное преобразование:
\[
\ln\bigl(1+e^{-2x^2}\bigr)=\ln\alpha^2-2x
\iff
F(x):=2x+\ln\bigl(1+e^{-2x^2}\bigr)=2\ln|\alpha|.
\]
\textbf{Симметрия.} В уравнение входит только $\alpha^2$, поэтому ответы при $\alpha$ и $-\alpha$ одинаковы (это видно и в таблице).

\subsection{Существование и единственность — прямое доказательство}
Хотя существование и единственность получатся из теоремы Банаха (раздел~4), полезно увидеть их «руками».
\[
F'(x)=2-\frac{4x\,e^{-2x^2}}{1+e^{-2x^2}}=2-\frac{4x}{e^{2x^2}+1}=2\bigl(1-G'(x)\bigr),
\]
где $G'(x)=\dfrac{2x}{e^{2x^2}+1}$ (см. ниже). Так как $|G'(x)|\le\frac12$ (лемма~1), имеем $F'(x)\ge1>0$:
$F$ строго возрастает. Кроме того, $2x\le F(x)\le 2x+\ln2$, поэтому $F(x)\to\pm\infty$ при $x\to\pm\infty$.
По теореме о промежуточном значении $F$ принимает значение $2\ln|\alpha|$ ровно один раз.

\textbf{Локализация корня.} Так как $0<\ln(1+e^{-2x^2})\le\ln2$, из $x^*=\ln|\alpha|-\frac12\ln(1+e^{-2x^{*2}})$ следует
\[
\ln|\alpha|-\tfrac12\ln2\;\le\;x^*\;<\;\ln|\alpha|.
\]
Это удобная проверка ответа: например, при $\alpha=10$ получаем $1.956\le x^*<2.3026$, и действительно $x^*\approx2.302573$
(при больших $|x|$ слагаемое $e^{-2x^2}$ ничтожно, и $x^*\approx\ln|\alpha|$).

\section{Метод простой итерации}
\subsection{Идея}
Уравнение записывают в виде $x=G(x)$ и строят последовательность $x_{n+1}=G(x_n)$ с начальным приближением $x_0$.
Если $G$ \emph{сжимает} расстояния, последовательность сходится к единственной неподвижной точке — решению.

\begin{theorem}[Банаха о сжимающем отображении, для $\mathbb R$]
Пусть $G\colon\mathbb R\to\mathbb R$ и существует $k\in[0,1)$ такое, что $|G(u)-G(v)|\le k|u-v|$ для всех $u,v$.
Тогда у $G$ есть единственная неподвижная точка $x^*$, и для любого $x_0$ последовательность $x_{n+1}=G(x_n)$ сходится к $x^*$, причём
\[
\text{(а) апостериорная оценка: } |x_n-x^*|\le\frac{|x_{n+1}-x_n|}{1-k};\qquad
\text{(б) априорная оценка: } |x_n-x^*|\le\frac{k^n}{1-k}|x_1-x_0|.
\]
\end{theorem}
\begin{proof}
По индукции $|x_{m+1}-x_m|\le k^m|x_1-x_0|$. Для $p\ge1$
\[
|x_{n+p}-x_n|\le\sum_{m=n}^{n+p-1}|x_{m+1}-x_m|\le|x_1-x_0|\sum_{m=n}^{\infty}k^m=\frac{k^n}{1-k}|x_1-x_0|\xrightarrow[n\to\infty]{}0,
\]
т.\,е. $(x_n)$ фундаментальна и (полнота $\mathbb R$) сходится к некоторому $x^*$. Сжатие непрерывно, поэтому, переходя к пределу в $x_{n+1}=G(x_n)$, получаем $x^*=G(x^*)$. Если $y^*=G(y^*)$ — ещё одна неподвижная точка, то $|x^*-y^*|=|G(x^*)-G(y^*)|\le k|x^*-y^*|$, откуда $x^*=y^*$.
Устремляя $p\to\infty$ в оценке выше, получаем (б). Оценка (а):
\[
|x_n-x^*|\le|x_n-x_{n+1}|+|x_{n+1}-x^*|=|x_n-x_{n+1}|+|G(x_n)-G(x^*)|\le|x_n-x_{n+1}|+k|x_n-x^*|,
\]
и остаётся перенести $k|x_n-x^*|$ влево и разделить на $1-k>0$.
\end{proof}

\subsection{Выбор отображения $G$}
Из логарифмической формы $2x+\ln(1+e^{-2x^2})=2\ln|\alpha|$ выражаем $x$:
\[
\boxed{\,G(x)=\ln|\alpha|-\frac12\ln\bigl(1+e^{-2x^2}\bigr)\,}
\]
Почему именно так? Нам нужно, чтобы $|G'|$ было меньше $1$ \emph{на всей прямой} и \emph{при всех} $\alpha$.
Здесь $\alpha$ входит только аддитивной константой $\ln|\alpha|$, которая исчезает при дифференцировании, а оставшаяся часть
$-\frac12\ln(1+e^{-2x^2})$ — ограниченная и очень «плоская» функция. Другие естественные варианты хуже:
\begin{itemize}
\item $x=x-\lambda\bigl(1+e^{-2x^2}-\alpha^2e^{-2x}\bigr)$: производная $1-\lambda(\dots)$ зависит от $\alpha$ и неограничена (из-за $e^{-2x}$), одного $\lambda$ на всю прямую не подобрать;
\item выражать $x$ из $e^{-2x^2}=\alpha^2e^{-2x}-1$: получаем $x^2=-\frac12\ln(\alpha^2e^{-2x}-1)$ — нужно следить за областью определения (аргумент логарифма должен быть положителен) и за знаком корня, удобной оценки нет.
\end{itemize}

\subsection{Доказательство сжимаемости}
\begin{lemma}
Для всех $x\in\mathbb R$ и всех $\alpha\ne0$: \quad $|G'(x)|\le\frac12$.
\end{lemma}
\begin{proof}
По цепному правилу $\bigl(\ln(1+e^{-2x^2})\bigr)'=\dfrac{-4x\,e^{-2x^2}}{1+e^{-2x^2}}$, поэтому
\[
G'(x)=-\frac12\cdot\frac{-4x\,e^{-2x^2}}{1+e^{-2x^2}}=\frac{2x\,e^{-2x^2}}{1+e^{-2x^2}}
=\frac{2x}{e^{2x^2}+1}
\]
(умножили числитель и знаменатель на $e^{2x^2}$). Так как $e^t\ge1+t$, при $t=2x^2$ знаменатель $e^{2x^2}+1\ge2+2x^2$. Значит,
\[
|G'(x)|\le\frac{2|x|}{2(1+x^2)}=\frac{|x|}{1+x^2}\le\frac12,
\]
где последнее неравенство — это $2|x|\le1+x^2\iff(|x|-1)^2\ge0$.
\end{proof}
По формуле Лагранжа $G(u)-G(v)=G'(\xi)(u-v)$ для некоторого $\xi$ между $u$ и $v$, поэтому
$|G(u)-G(v)|\le\frac12|u-v|$: $G$ — сжатие с константой $k=\frac12$.

\textbf{Замечание (точная константа).} Численно $\sup_x|G'(x)|\approx0.39295$ (достигается при $|x|\approx0.6078$).
Мы используем доказуемую оценку $k=\frac12$: она немного грубее, зато строго обоснована; от неё зависит только критерий остановки.

\subsection{Критерий остановки}
По апостериорной оценке (а) с $k=\frac12$:
\[
\frac{|x_{n+1}-x_n|}{1-k}\le\varepsilon\iff|x_{n+1}-x_n|\le\frac{\varepsilon}{2}=5\cdot10^{-13}\quad\Longrightarrow\quad|x_n-x^*|\le\varepsilon=10^{-12}.
\]
Поэтому в ответ выдаётся $x_n$ — приближение, для которого \emph{гарантирована} точность $\varepsilon$, а $n$ — его номер (как в образце).

\textbf{Априорная оценка числа шагов.} $x_0=0$, $x_1=\ln|\alpha|-\frac12\ln2$. Из (б) точность $\varepsilon$ гарантирована, как только
$\frac{k^n}{1-k}|x_1-x_0|\le\varepsilon$, т.\,е. $n\ge\log_2\frac{2|x_1|}{\varepsilon}$. Например, при $\alpha=2$: $|x_1|=\frac12\ln2\approx0.347$,
$n\ge\log_2(6.9\cdot10^{11})\approx39.3$. Фактически хватает $26$ шагов (критерий (а) точнее, чем (б)).

\subsection{Почему число итераций разное}
Вблизи корня ошибка убывает примерно как $|x_{n+1}-x^*|\approx|G'(x^*)|\,|x_n-x^*|$, т.\,е. реальная скорость определяется числом $|G'(x^*)|$, а не грубой константой $k$:
\begin{center}
\begin{tabular}{c|c|c|c}
$|\alpha|$ & $x^*$ & $G'(x^*)$ & $n$\\\hline
1 & $-0.302822$ & $-0.275$ & 21\\
2 & $0.430809$ & $0.352$ & 26\\
3 & $1.045342$ & $0.211$ & 19\\
4 & $1.375027$ & $0.061$ & 12\\
6 & $1.790942$ & $0.0059$ & 7\\
10 & $2.302573$ & $0.00011$ & 5
\end{tabular}
\end{center}
При больших $|\alpha|$ корень далеко от нуля, $G'(x^*)\approx 2x^*e^{-2x^{*2}}$ очень мал — сходимость почти мгновенная.
Самый медленный случай — $|\alpha|$ около $2$, когда $x^*$ близок к точке максимума $|G'|$.

\subsection{Вычислительные тонкости}
\begin{itemize}
\item $\ln(1+t)$ при малом $t=e^{-2x^2}$ вычисляется функцией \texttt{std::log1p(t)}: обычный \texttt{log(1+t)} при $t<10^{-16}$ дал бы $\log 1=0$ с потерей всех значащих цифр $t$.
\item При больших $|x|$ величина $e^{-2x^2}$ становится машинным нулём (\emph{underflow}) — это безопасно: тогда $G(x)=\ln|\alpha|$, что верно с машинной точностью.
\item $\ln|\alpha|$ вычисляется \emph{один раз} до цикла.
\item Критерий $|x_{n+1}-x_n|\le5\cdot10^{-13}$ достижим в арифметике \texttt{double}: ошибка округления при вычислении $G$ порядка $10^{-16}\max(1,|x|)$, а для любых конечных \texttt{double} $\alpha\ne0$ корень удовлетворяет $|x^*|\le|\ln|\alpha||+1<746$. На всякий случай в коде есть страховочный лимит \texttt{MAX\_ITER}.
\end{itemize}

\section{Разбор кода}
\VerbatimInput[numbers=left,fontsize=\small,frame=single,numbersep=4pt]{iterative_solver.cpp}

\begin{description}
\item[Строки 1–4.] Подключаем \texttt{cmath} (\texttt{exp}, \texttt{log}, \texttt{log1p}, \texttt{abs}, \texttt{isfinite}), \texttt{cstdlib} (\texttt{strtod}, \texttt{strtol} — разбор аргументов), \texttt{iomanip} (форматирование таблицы), \texttt{iostream} (вывод).
\item[Строки 11–13.] Константы: требуемая точность $\varepsilon=10^{-12}$, константа сжатия $k=\frac12$ (лемма~1) и страховочный лимит итераций.
\item[Строки 15–17.] Функция \texttt{G} — отображение $G(x)=\ln|\alpha|-\frac12\ln(1+e^{-2x^2})$. Параметр передаётся уже в виде $\ln|\alpha|$.
\item[Строки 19–27.] Функция \texttt{solve(alpha)} сначала печатает $\alpha$; при $\alpha=0$ сообщает, что решений нет (раздел 2.1).
\item[Строки 29–31.] Вычисляем $\ln|\alpha|$, начальное приближение $x_0=0$, счётчик $n=0$.
\item[Строки 33–48.] Цикл метода простой итерации. На каждом шаге \texttt{new\_x} $=x_{n+1}=G(x_n)$, \texttt{step} $=|x_{n+1}-x_n|$.
Если $\texttt{step}/(1-k)\le\varepsilon$, то по апостериорной оценке $|x_n-x^*|\le\varepsilon$ — выходим, \emph{не} заменяя $x$ на \texttt{new\_x}, так что в \texttt{x} лежит $x_n$, а в \texttt{n} — его номер. Иначе переходим к следующему приближению.
\item[Строки 50–52.] Печать $x_n$ с 15 знаками после точки и номера $n$.
\item[Строки 55–59.] Заголовок таблицы.
\item[Функция \texttt{main}.] Два режима, как в образце.
\begin{itemize}
\item Один аргумент — вещественное $\alpha$. \texttt{strtod} разбирает строку; проверяем, что строка — целиком число (\texttt{end} указывает на конец строки) и что оно конечно (\texttt{isfinite}).
\item Два аргумента — целые $a<b$; решаем уравнение для всех целых $\alpha\in[a,b]$, включая концы.
\item Иначе — подсказка по использованию и код возврата 1.
\end{itemize}
\end{description}

\textbf{Отличия от образца.} Образец решает \emph{систему} из двух уравнений (вектор $(x,y)$, норма $\|\cdot\|_1$, $k=\frac34$); у нас \emph{одно} уравнение, поэтому одна переменная $x$, модуль вместо нормы и $k=\frac12$. Добавлена обработка $\alpha=0$ (у образца решение есть при всех $\alpha$, у нас — нет) и страховочный лимит итераций. Структура программы, формат вывода и интерфейс командной строки — те же.

\section{Проверка результатов}
Корни были независимо перепроверены в Python с библиотекой \texttt{mpmath} (40 значащих цифр, метод Ньютона \texttt{findroot}) для всех целых $\alpha$ от $-10$ до $10$, $\alpha\ne0$.
Фактическая ошибка $|x_n-x^*|$ всюду не превосходит $4.5\cdot10^{-13}<10^{-12}$ (наибольшая — при $|\alpha|=8$), а невязка
$1+e^{-2x_n^2}-\alpha^2e^{-2x_n}$ по модулю не больше $1.4\cdot10^{-12}$. Значения $\pm\alpha$ дают одинаковые ответы, все корни лежат в отрезке
$[\ln|\alpha|-\frac12\ln2,\ \ln|\alpha|)$ из раздела~2.3.

Дополнительно проверены крайние случаи: $\alpha=10^{300}$ ($x\approx690.7755$, 2 итерации), $\alpha=10^{-300}$ ($x\approx-690.7755$), $\alpha=0.5$ ($x\approx-0.81175$, 27 итераций), некорректный ввод (\texttt{abc}) — выдаётся сообщение об ошибке.

\section{Возможные вопросы на сдаче}
\begin{qa}
\textbf{Почему метод сходится при любом $\alpha$?} — Производная $G'$ не зависит от $\alpha$ и по модулю $\le\frac12$ на всей прямой; по теореме Банаха итерации сходятся из любой начальной точки.

\textbf{Почему решение единственно?} — Неподвижная точка сжатия единственна (или: $F'(x)\ge1>0$, функция $F$ строго возрастает).

\textbf{Что будет при $\alpha=0$?} — Решений нет: левая часть $\ge1$, правая равна $0$.

\textbf{Откуда критерий остановки?} — Из апостериорной оценки $|x_n-x^*|\le\frac{|x_{n+1}-x_n|}{1-k}$; при $k=\frac12$ это $|x_{n+1}-x_n|\le5\cdot10^{-13}$.

\textbf{Почему печатается $x_n$, а не $x_{n+1}$?} — Именно для $x_n$ гарантирована оценка $\le\varepsilon$ (так же сделано в образце). На самом деле $x_{n+1}$ ещё точнее: $|x_{n+1}-x^*|\le\frac{k}{1-k}|x_{n+1}-x_n|$.

\textbf{Почему при $|\alpha|=2$ итераций больше всего?} — Скорость определяется $|G'(x^*)|$; при $|\alpha|\approx2$ корень $x^*\approx0.43$ близок к точке максимума $|G'|$ ($\approx0.61$), и $|G'(x^*)|\approx0.35$ — наибольшее среди целых $\alpha$.

\textbf{Зачем \texttt{log1p}?} — Для точного вычисления $\ln(1+t)$ при очень малых $t$.

\textbf{Можно ли взять $k\approx0.393$?} — Да, это точная верхняя грань $|G'|$, но её надо обосновать (найдена численно); $k=\frac12$ доказывается в одну строку, а на ответ и гарантию точности это не влияет.
\end{qa}
\end{document}
